\section{Comparison with the Physical Atmosphere}
\label{sec:physical-atmosphere}

The inverse construction of Sec.~\ref{sec:inverse-reconstruction} shows that a
stratified refractive medium can be designed mathematically to reproduce a
prescribed horizon law above a flat opaque boundary.  The remaining issue is
whether the required profile resembles the refractive-index structure of an
ordinary atmosphere.  The decisive difference is the sign of the vertical
refractive-index gradient.

\subsection{Gradient required by the optical construction}

The Earth-scale exponential model is
\begin{equation}
    n_{\mathrm{e}}(h)=n_0 e^{h/R},
    \label{eq:required-exponential-profile}
\end{equation}
which gives
\begin{equation}
    \frac{dn_{\mathrm{e}}}{dh}
    =
    \frac{n_{\mathrm{e}}(h)}{R}
    >0.
    \label{eq:required-positive-gradient}
\end{equation}
The exact profile reconstructed from the spherical horizon law,
Eq.~\eqref{eq:exact-spherical-reconstruction-profile}, has the same qualitative
behaviour.  Indeed, writing
\begin{equation}
    \frac{n_{\mathrm{s}}^2(h)}{n_0^2}
    =
    1+
    \frac{(R+h)^2h(2R+h)}{R^4},
\end{equation}
the polynomial multiplying $R^{-4}$ is strictly increasing for $h>0$.
Therefore,
\begin{equation}
    \frac{dn_{\mathrm{s}}}{dh}>0
    \qquad (h>0).
    \label{eq:reconstructed-positive-gradient}
\end{equation}
Both the asymptotically matched exponential profile and the exact inverse
profile thus require the refractive index to increase with altitude.

This sign also agrees with the ray geometry derived above.  In a stratified
isotropic medium, rays bend toward regions of larger refractive index.  A
profile with $dn/dh>0$ therefore bends rays upward, allowing a ray that joins
two elevated endpoints to pass through a lower turning point and graze the
opaque boundary between them.  This turning point is the optical analogue of a
geometrical tangency point on a convex surface.

\subsection{Ordinary atmospheric stratification}

Except under strong local inversions, the refractive index of air decreases
with altitude as density and pressure decrease
\cite{ciddor1996refractive,ussa1976,bean1966radio,itur2019p453}.  A simple
idealised profile is
\begin{equation}
    n_{\mathrm{a}}(h)
    \simeq
    1+(n_0-1)e^{-h/H},
    \label{eq:ordinary-atmosphere-profile}
\end{equation}
where $H$ is an atmospheric scale height.  Its derivative is
\begin{equation}
    \frac{dn_{\mathrm{a}}}{dh}
    =
    -\frac{n_0-1}{H}e^{-h/H}
    <0.
    \label{eq:ordinary-negative-gradient}
\end{equation}
The ordinary atmosphere therefore has the opposite gradient sign to that in
Eqs.~\eqref{eq:required-positive-gradient} and
\eqref{eq:reconstructed-positive-gradient}.  Rays bend toward the denser,
higher-index air below rather than upward toward regions of increasing index.

This mismatch is qualitative rather than merely quantitative.  Changing the
scale height or the small surface refractivity $n_0-1$ modifies the magnitude
of the gradient but not its sign.  A monotonically decreasing profile cannot
be continuously adjusted into the required increasing profile without passing
through a qualitatively different stratification.

\subsection{Scale of the required variation}

For the exponential construction, the fractional gradient is
\begin{equation}
    \frac{1}{n}\frac{dn}{dh}
    =
    \frac{1}{R}.
    \label{eq:required-fractional-gradient}
\end{equation}
Near the boundary,
\begin{equation}
    n(h)-n_0
    \simeq
    n_0\frac{h}{R}.
    \label{eq:required-near-surface-index-change}
\end{equation}
The exact inverse profile has the same linear term, as shown in
Eq.~\eqref{eq:spherical-profile-small-height}.  The leading requirement is
therefore not an arbitrarily chosen local anomaly; it is fixed by the radius
$R$ of the spherical horizon law being reproduced.

The absolute change over terrestrial heights is small because $h/R\ll1$.
Its small magnitude, however, does not remove the sign conflict.  Weak positive
and negative gradients bend rays in opposite directions and therefore produce
different turning-point structures.

\subsection{Exceptional atmospheric conditions}

Real atmospheres need not be perfectly monotonic.  Temperature inversions and
sharp boundary layers can produce anomalous refraction, ducting, mirages, and
multiple ray paths \cite{bean1966radio,itur2019p453,nener2003analytical}.
Such effects may substantially shift an apparent horizon
or distort individual objects.  They do not, however, render the ordinary
atmosphere globally equivalent to the smooth profile reconstructed above.
Exact reproduction of the spherical visibility law for arbitrary endpoint
heights requires the prescribed positive profile throughout the relevant
height interval, rather than a local layer that occasionally bends selected
rays in the required direction.

Two distinct statements should therefore be separated.  First, inverse
geometrical optics shows that a mathematical refractive profile above a flat
boundary can reproduce the same limiting horizon law as a sphere.  Second, the
profile required for this equivalence is not that of an ordinary atmosphere,
because its vertical gradient has the opposite sign.  The construction is
therefore an exact optical equivalence, but not a model of standard atmospheric
refraction.
